\section{Vacuity of the along-trajectory residual}
\label{app:vacuity}
% =====================================================================

Physics-informed training conventionally penalizes the PDE residual evaluated on
the model's own trajectory. The following makes precise why that penalty is
uninformative for the present architecture, and motivates the steady-state
formulation adopted instead.

\begin{proposition}[The trajectory residual is an integrator diagnostic]
\label{prop:vacuity}
Define $R_{\Delta t}(\vz;\theta) := \Delta t^{-1}\big(S_{\Delta t}\vz - \vz\big)
- \big(\mathcal{A}_{\mD}\vz + f_\theta(\vz)\big)$. Then for every
$\theta$ in a bounded parameter set and every sufficiently regular $\vz$,
\begin{equation}
\|R_{\Delta t}(\vz;\theta)\| \;=\; \Delta t \cdot
\Big\| \tfrac{1}{2}\big[\mathcal{A}_{\mD},\,\mathcal{F}_\theta\big]\vz \Big\|
\;+\; O(\Delta t^2),
\label{eq:residual}
\end{equation}
where $\mathcal{F}_\theta \vz := f_\theta(\vz)$ and $[\cdot,\cdot]$ is the
commutator. In particular $R_{\Delta t}\to0$ as $\Delta t\to0$ \emph{uniformly in
$\theta$}, and $R_{\Delta t}$ does not depend on the observed data.
\end{proposition}

\begin{proof}[Proof sketch]
Expanding $S_{\Delta t}$ by BCH and subtracting the exact generator leaves the
leading commutator term at order $\Delta t$; all data-dependent quantities cancel
because $S_{\Delta t}$ and the generator are evaluated at the same state $\vz$.
\end{proof}

Two consequences follow. First, minimizing
$\|R_{\Delta t}\|^2$ cannot fit the model to observations, since the observations
do not appear in \eqref{eq:residual}; it drives $\theta$ toward operators whose
diffusive and reactive parts commute, an inductive bias unrelated to the data.
Second, the gradient signal it supplies is $O(\Delta t)$ and therefore vanishes
in the very limit in which the discretization is most faithful. The empirical
scaling is confirmed in Table~\ref{tab:theory} (observed order $1.00$).

The constraint retained instead is the \emph{steady-state} residual
\eqref{eq:steady}, which asks that the measured configuration be a fixed point of
the learned operator. This does involve the data, through $\vz_T$, and it
constrains $\mD_\theta$ and $f_\theta$ jointly: the set of pairs satisfying it is
a proper subset of parameter space, characterized next.

% =====================================================================
